% PARDON-PRIMITIVE1
\section*{PARDON-PRIMITIVE1: colimit quotient and primitive continuation structure}

\textbf{STATUS: COMPUTATIONAL finite experiment plus KNOWN-THEOREM literature mapping. P versus NP remains OPEN.}

\paragraph{Literature audit.}
Pardon's 2023 curve-counting result shows that, in its geometric setting, global enumerative
invariants are determined by local-curve values. Behrend's 2026 zero-cycle analogue makes
the algebraic architecture explicit. Pardon homology is first defined as a colimit over
families, with pullback-related representatives identified. Only after these relations are
imposed does the product/coproduct make the colimit into a graded commutative and
cocommutative Hopf algebra. Multiplicative theories are group-like, primitive theories are
described as connected theories, and Proposition~2.14 gives an exponential/logarithm
bijection between multiplicative and primitive theories.

This ordering is important for the computational analogy:
\[
\boxed{
\text{quotient certified equivalent realizations first}
\quad\longrightarrow\quad
\text{extract primitive/connected structure second}.
}
\]
Primitive extraction alone cannot make wrappers disappear.

Schmitt's incidence Hopf algebras and the graph Hopf algebra of Humpert--Martin provide a
more direct combinatorial precedent: disjoint union supplies the product and restriction/
splitting the coproduct, while primitive structure isolates connected information.

\paragraph{Finite computational analogue.}
For a promise classifier and a fixed source-coordinate split, form the binary continuation
matrix $M$ whose one-entries are legal UNSAT completions. Define
\[
\rho(M)=\operatorname{rank}_{\mathbb F_2}M,
\qquad
\kappa(M)=\log_2\rho(M).
\]
This is not a Pardon primitive element. It is a deliberately modest scalar analogue:
Kronecker products obey
\[
\rho(M\otimes N)=\rho(M)\rho(N),
\]
so
\[
\boxed{\kappa(M\otimes N)=\kappa(M)+\kappa(N)}.
\]
Thus logarithm converts multiplicative independent-product inflation into an additive quantity.

\paragraph{Wrapper test.}
On the natural $16|16$ clause split of $D(4,8)$, the balanced $k=4$ continuation block has
shape $84\times84$, 112 one-entries, and
\[\rho=56.\]
It has exactly 56 distinct nonzero row signatures. Duplicating every row (COPY), appending
zero rows (DEAD), or duplicating rows to model the semantic effect of a tautological fork all
leave the rank exactly 56. Hence
\[\kappa=\log_2 56\approx5.80735.\]
The qualification ``semantic'' is load-bearing: this happens because the continuation
relation has already forgotten syntax. A raw syntactic DAG algebra would retain a wrapper
unless an explicit quotient relation identifies it with the unit.

\paragraph{Independent-product test.}
The $D(4,8)$ $k=2$ block has rank 16. Its explicitly constructed Kronecker square has rank
256, so
\[\kappa(M)=4,\qquad \kappa(M\otimes M)=8.\]

\paragraph{Hostile controls.}
For a $6|6$ split of parity, the continuation matrix over $\mathbb F_2$ has rank 2 and
$\kappa=1$, despite parity having arbitrarily large ordinary decision depth. For six-bit
equality on each side, the matrix is the $64\times64$ identity and
$\rho=64,\ \kappa=6$. Thus the score measures cross-cut semantic distinction, not raw
decision count.

\paragraph{Primitive decomposition of the fixed-eight control.}
The natural $D(4,8)$ balanced $k=4$ support graph has 56 connected components and every
component is complete bipartite:
\[
16K_{2,1}\sqcup16K_{1,1}\sqcup16K_{1,2}\sqcup4K_{1,4}\sqcup4K_{4,1}.
\]
Consequently
\[
\operatorname{rank}_{\mathbb F_2}M
=\#\{\text{nonzero row types}\}
=\#\{\text{support components}\}=56.
\]

The same phenomenon persists for the natural $40|40$ split of $D(5,8)$. The nonempty
weight blocks have ranks $1,28,144,36,1$ for $k=0,2,4,6,8$, and every connected
support component in every tested block is complete bipartite. At $k=4$ there are 144
components; the largest tested component has only 22 support vertices.

\paragraph{Interpretation.}
This is a negative result for using a Pardon-style primitive layer as a fixed-$m=8$ hardness
certificate. The global continuation relation decomposes into many small local biclique pieces
rather than one growing irreducible component. This matches the known polynomial conflict
representation of $D(n,8)$.

The useful architecture is instead
\[
\text{raw realizations}
\to
\text{certified semantic/refinement quotient}
\to
\text{independent-composition product}
\to
\text{continuation relation}
\to
\text{multiplicative invariant}
\to
\text{primitive/log score}.
\]
The quotient step addresses wrappers; the logarithm step addresses independent-product inflation.
Neither step supplies an arbitrary-algorithm SAT lower bound.

\paragraph{Next target.}
Test the same connected-support and log-rank quantities for $m=\alpha n$ around the random
3-SAT hard region. If the continuation relation again decomposes into uniformly small easy
components, retire this primitive-decomposition route as a hardness invariant. If large
connected primitive components emerge, compare them against FBDD/branching-program state,
resolution width, and affine/XOR hostile controls.

\begin{thebibliography}{9}
\bibitem{Pardon23} J.~Pardon, \emph{Universally counting curves in Calabi--Yau threefolds}, arXiv:2308.02948, 2023.
\bibitem{Behrend26} K.~Behrend, \emph{A Pardon Algebra for Zero-cycles}, arXiv:2604.05534, 2026.
\bibitem{Schmitt94} W.~R.~Schmitt, \emph{Incidence Hopf Algebras}, J. Pure Appl. Algebra 96 (1994), 299--330.
\bibitem{HumpertMartin12} B.~Humpert and J.~L.~Martin, \emph{The Incidence Hopf Algebra of Graphs}, SIAM J. Discrete Math. 26 (2012), 555--570.
\end{thebibliography}

% END_OF_NOTE::PARDON-PRIMITIVE1